% !TeX root = ../抽象代数笔记.tex
% 由主文档通过 \input 载入；请编译主文档。
\setcounter{chapter}{15} % 保留原书章号。
\chapter{分裂域、可分扩张与正规扩张}\label{ch:splitting}

\begin{description}
  \item[学习目标] 理解包含全部根、根无重复和嵌入数量之间的联系，准确掌握伽罗瓦扩张的条件。
  \item[先修知识] 第\ref{ch:extensions}章、多项式不可约性与形式导数。
\end{description}

\section{分裂域的构造与唯一性}
对域 $F$ 上的非恒定多项式 $f$，定义分裂域（splitting field）为 $f$ 分裂且由其全部根生成的扩域。通过逐次取不可约因子并添加根证明存在性，通过扩张嵌入证明分裂域在 $F$-同构意义下唯一；唯一性不表示存在唯一一个同构。

\section{可分多项式与可分扩张}
定义可分多项式（separable polynomial）为在分裂域中没有重根的多项式，代数元素可分指其最小多项式可分。代数扩张可分指其所有元素可分。证明特征为 $0$ 的域上不可约多项式可分，并用 $\mathbb F_p(t)$ 上的 $x^p-t$ 展示正特征下不可分现象，分别验证不可约性与重根性质。

\section{正规扩张与伽罗瓦扩张}
定义代数扩张 $K/F$ 正规（normal），是指 $F[x]$ 中每个在 $K$ 有根的不可约多项式都在 $K$ 完全分裂。证明有限扩张正规当且仅当它是某个多项式的分裂域。定义伽罗瓦扩张（Galois extension）为正规且可分的代数扩张；本书随后主要处理有限情形。

\section{嵌入数量与本原元}
定义固定基域 $F$ 的域嵌入，证明有限扩张 $K/F$ 到一个代数闭包中的 $F$-嵌入数不超过 $[K:F]$，等号成立当且仅当扩张可分。证明本原元定理（primitive element theorem）：有限可分扩张可由单个元素生成；说明这是一条带条件的存在性结论，不能推广为任意有限扩张都为单扩张。

\section{代数闭包与例题练习}
定义代数闭包（algebraic closure）；其一般存在性与唯一性证明涉及额外集合论工具，列为选学并注明所用原则。核心例题比较 $\mathbb Q(\sqrt[3]{2})/\mathbb Q$ 与 $x^3-2$ 的分裂域。以 $\mathbb Q$ 在 $\mathbb C$ 中的代数闭包为代数而非有限的例子，通过任意高次数的艾森斯坦多项式说明次数无界。

